\section{Conclusiones}

En esta primera entrega se ha establecido la estructura base del proyecto. A partir de la identificación de una problemática real y frecuente en entornos organizacionales (la ineficiencia en la coordinación grupal) se definió una solución centrada en un agente con un LLM como núcleo. El estudio de factibilidad evidencia que el proyecto es técnica y económicamente realizable dentro del plazo de un semestre, con un stack tecnológico relativamente accesible. La adopción de Scrum como metodología de desarrollo, junto con la definición clara de roles y un diseño arquitectural modular, sienta las condiciones para un avance iterativo y controlado hacia el MVP. Las siguientes etapas del proyecto se enfocarán en la implementación progresiva de los sprints definidos, con miras a demostrar en la Evaluación 2 un avance funcional de al menos el 50\% de los objetivos comprometidos, tal como se pide.
